\documentclass[11pt]{article}
\usepackage{fontspec}
\setmainfont{Noto Serif CJK SC}[BoldFont={Noto Serif CJK SC SemiBold}]
\setsansfont{Noto Sans CJK SC}
\XeTeXlinebreaklocale "zh"
\XeTeXlinebreakskip = 0pt plus 1pt
\usepackage[a4paper,margin=2.3cm]{geometry}
\usepackage{amsmath,amssymb}
\usepackage{booktabs,array}
\usepackage[hidelinks]{hyperref}
\linespread{1.25}
\setlength{\parskip}{0.45em}
\setlength{\parindent}{0pt}
\newcommand{\cl}{c_{\text{late}}}
\newcommand{\cs}{c_{\text{slack}}}
\newcommand{\1}{\mathbf 1}
\newcounter{prop}
\newenvironment{prop}[1][]{\refstepcounter{prop}\par\medskip\noindent\textbf{命题 \theprop}\if\relax\detokenize{#1}\relax\else\ \textbf{（#1）}\fi\textbf{.}\ }{\par\medskip}
\newenvironment{pf}{\par\noindent\textbf{证明.}\ }{\par\noindent\hfill QED\par\medskip}

\begin{document}

{\Large\textbf{为什么用分位数设定预计日期，以及校准的保证到哪里为止}}

\medskip
这份说明回答两个问题。第一，为什么预计日期取配送天数的分位数，而不是平均值。第二，校准后的“覆盖率达标”到底指什么。我会区分三个经常被混在一起的量：校准集上的比例、可交换性下的概率，以及评价期的实际比例。这里用到的都是标准结果：分位数是不对称线性损失的最小点，以及分割保形校准背后的排名论证。我按本项目的形式写出并证明它们，再用到项目上；它们都不是新结论。

\bigskip
\textbf{1.\ 不对称损失为什么对应分位数}

设实际配送天数为整数 \(Y\ge1\)，\(\mathbb E[Y]<\infty\)；预计天数为整数 \(d\ge1\)。日期太短造成延误，日期太长留下空等。我用线性损失表示这两种代价：
\[
L(d,Y)=\cl\,(Y-d)_+ +\cs\,(d-Y)_+,\qquad \cl,\cs>0 .
\]
这是代价随天数线性增加的简化假设。记 \(R(d)=\mathbb E[L(d,Y)]\)，\(F(d)=P(Y\le d)\)，\(q=\cl/(\cl+\cs)\)。

\begin{prop}
\(d^\ast=\min\{d\ge1:\ F(d)\ge q\}\) 使 \(R(d)\) 最小。若 \(F(d^\ast)=q\)，则 \(d^\ast+1\) 同样最优。
\end{prop}

\begin{pf}
比较承诺 \(d\) 天与 \(d+1\) 天。逐单来看，
\[
L(d+1,Y)-L(d,Y)=
\begin{cases}
-\cl, & Y>d,\\
+\cs, & Y\le d .
\end{cases}
\]
取期望，
\[
\Delta(d):=R(d+1)-R(d)=\cs F(d)-\cl\bigl(1-F(d)\bigr)=(\cl+\cs)\bigl(F(d)-q\bigr).
\]
\(F\) 单调不减，所以 \(\Delta\) 单调不减。对 \(d<d^\ast\)，\(F(d)<q\)，故 \(\Delta(d)<0\)；对 \(d\ge d^\ast\)，\(\Delta(d)\ge0\)。因此 \(R\) 在 \(d^\ast\) 之前严格下降，之后不再下降，\(d^\ast\) 是最小点。又 \(F(0)=0<q\)，所以 \(d^\ast\ge1\)。若 \(F(d^\ast)=q\)，则 \(\Delta(d^\ast)=0\)，即 \(R(d^\ast+1)=R(d^\ast)\)。
\end{pf}

对单笔订单，用下单时的信息 \(X=x\)，把 \(F\) 换成 \(F(d\mid x)=P(Y\le d\mid X=x)\)，最优日期就是条件 \(q\) 分位数。取 \(\cl=q\)、\(\cs=1-q\) 时，\(L\) 就是分位数（pinball）损失。

这个损失说明了分位数目标的决策含义，但它不等于每个模型的拟合目标。项目中的线性模型在 \(\log(1+Y)\) 尺度上拟合分位数回归，再换回天数；HGB 直接在天数上拟合；模型选择和评价都回到天数尺度。分位数在单调变换下不变，所以两种尺度瞄准的是同一个条件分位数，但有限样本中的拟合目标并不相同。

平均值最小化的是平方损失，它把晚到和空等看成对称的代价。一般而言，平均值并不最小化这里定义的不对称损失；在某些分布下，它可能恰好与最优分位数相同。选择 \(q=0.90\) 等价于假设 \(\cl/\cs=9\)，即晚到一天的代价是多留一天的九倍。这是我选定的目标，不是从 Olist 数据中估计出的真实成本比。

\bigskip
\textbf{2.\ 校准：三个不同的“覆盖率”}

模型给出的 \(f(x)\) 只是估计值，直接使用未必达到目标。因此用独立留出的校准集修正。设 \(f\) 已拟合完毕，且没有用到校准集。在 \(n\) 个校准订单上计算
\[
r_i=Y_i-f(X_i),\qquad k=\lceil (n+1)q\rceil\le n,\qquad s=r_{(k)},\qquad U(x)=f(x)+s ,
\]
其中 \(r_{(k)}\) 是第 \(k\) 小的残差。注意 \(Y_i\le U(X_i)\iff r_i\le s\)。

\textbf{第一个量：校准集上的比例，至少 \(k/n\)。}

\begin{prop}
对任意数据，\(\displaystyle \hat c_{\mathrm{cal}}:=\frac1n\sum_{i=1}^n \1\{Y_i\le U(X_i)\}\ \ge\ \frac kn\)。
\end{prop}

\begin{pf}
按定义，至少有 \(k\) 个 \(r_i\) 不超过 \(r_{(k)}=s\)（有并列时可能更多）。由 \(Y_i\le U(X_i)\iff r_i\le s\) 即得。
\end{pf}

这个不等式不依赖任何概率假设，即使未来配送状况已经改变，它仍然成立。该不等式是按指定校准规则计算后的必然结果；仅满足它，既不能证明实现正确，也不能保证未来订单的覆盖率。例如把平移量错设成 1,000 天，同样会满足它。更直接的实现检查，是独立重算第 \(k\) 个残差，再与冻结的平移量核对。

\textbf{第二个量：新订单被覆盖的概率，至少 \(k/(n+1)\)。}

\begin{prop}
若 \((X_1,Y_1),\dots,(X_n,Y_n),(X_{n+1},Y_{n+1})\) 可交换，则
\[
P\bigl(Y_{n+1}\le U(X_{n+1})\bigr)\ \ge\ \frac{k}{n+1}\ \ge\ q .
\]
\end{prop}

\begin{pf}
因为 \(f\) 固定，\(r_1,\dots,r_{n+1}\) 也可交换。随机打破并列后，把 \(n+1\) 个残差排序，记新订单残差 \(r_{n+1}\) 的名次为 \(\rho\)。由可交换性，\(\rho\) 在 \(\{1,\dots,n+1\}\) 上均匀分布，所以 \(P(\rho\le k)=k/(n+1)\)。若 \(\rho\le k\)，则前 \(n\) 个残差中排在 \(r_{n+1}\) 之前的至多 \(k-1\) 个，故其中第 \(k\) 小者不小于 \(r_{n+1}\)，即 \(r_{n+1}\le s\)。因此 \(P(r_{n+1}\le s)\ge k/(n+1)\)。最后，\(k\ge(n+1)q\)。
\end{pf}

分母是 \(n+1\)，因为新订单本身也是这 \(n+1\) 个可交换样本之一。这个概率是对“抽到哪一份校准集”和“来的是哪一笔新订单”同时取平均。它不保证每个月、每个州都达标，也不保证某一份已观察到的校准集之后的覆盖概率达标。随机打破并列只用于证明，不改变实际算法。

\textbf{第三个量：评价期的实际比例。}

设评价期有 \(m\) 笔订单，
\[
\hat c_{\mathrm{eval}}:=\frac1m\sum_{j=1}^m \1\{Y_j\le U(X_j)\}.
\]
它是一个观察到的比例，不是概率。如果每笔评价订单都与校准订单可交换，由命题 3 得 \(\mathbb E[\hat c_{\mathrm{eval}}]\ge k/(n+1)\)；即使如此，有限的 \(m\) 也会带来抽样波动，单次观察可能低于或高于这个值。若可交换性不成立，它只描述这批评价订单本身。

\begin{center}
\small
\begin{tabular}{>{\raggedright\arraybackslash}p{0.25\textwidth}>{\raggedright\arraybackslash}p{0.33\textwidth}>{\raggedright\arraybackslash}p{0.33\textwidth}}
\toprule
量 & 能说明什么 & 不能说明什么\\
\midrule
\(\hat c_{\mathrm{cal}}\ge k/n\)（命题 2） & 按指定规则计算后的必然结果 & 实现正确；未来订单也会达到目标\\
\(P(Y_{n+1}\le U)\ge k/(n+1)\)（命题 3） & 在可交换性下，对校准集和新订单取平均的覆盖至少为 \(q\) & 每份校准集、每个州、每个月都达标；C 与 E 的可交换性并未确立，因此不对评价期断言这一保证\\
\(\hat c_{\mathrm{eval}}\) & 冻结策略在这批评价订单上的经验表现 & 全部订单或未来运营中的可靠性\\
\bottomrule
\end{tabular}
\end{center}

一个简单反例说明第一个量与第三个量可以完全脱节。模型始终预测 10 天，校准订单也全部在 10 天送达，于是 \(s=0\)，\(\hat c_{\mathrm{cal}}=100\%\)。若未来订单都要 15 天，\(\hat c_{\mathrm{eval}}=0\)。反过来，若未来订单都只要 5 天，\(\hat c_{\mathrm{eval}}=100\%\)，但每单多承诺 5 天。两种情况下命题 2 都成立；命题 3 的可交换性却不成立，所以不能沿用其保证（\href{https://arxiv.org/abs/2202.13415}{Barber 等人}）。

\bigskip
\textbf{3.\ 整数日期}

预计日期必须是整数。因为 \(Y\) 是整数，\(Y\le U\iff Y\le\lfloor U\rfloor\)，所以向下取整与连续边界覆盖完全相同的订单；最低 1 天的限制只会扩大覆盖。项目登记的规则是 \(\max\{1,\lceil U\rceil\}\)。它覆盖的订单包含 \(U\) 覆盖的订单，所以命题 2、3 中的下界对它同样成立。把 \(U\) 向上取整为 \(\lceil U\rceil\)，增加的天数小于一天。最低 1 天的限制是另一项调整，只会提高覆盖。

\bigskip
\textbf{4.\ 回到 Olist}

校准集有 \(n=17{,}056\) 笔订单，\(k=15{,}352\)，所以 \(k/n\approx90.01\%\)，\(k/(n+1)\approx90.00\%\)。以线性模型为例，三个量分别是：
\begin{itemize}\setlength{\itemsep}{0pt}
\item 校准集上：登记的向上取整日期覆盖 91.19\%，不低于 \(k/n\)，与命题 2 一致；独立重算第 \(k\) 个残差得到的平移量（6.1479 天）与冻结值相同；
\item 可交换性下的保证：至少 \(k/(n+1)\)，但 C（2017 年 10--12 月）与评价期 E（2018 年 4--7 月）是否可交换，本身就是问题；
\item 评价期上：19,950 笔订单的实际覆盖为 96.18\%。
\end{itemize}

E 上的过覆盖提示，校准期与评价期之间的关系值得进一步检查。这个现象与配送时间变化、样本筛选及订单构成变化都可能相容。本说明没有做专门的可交换性检验，也不能分离这些因素的贡献。因此，我把它作为诊断线索，而不是对某一个原因的确认。此外，全部结论只适用于最终签收的单卖家订单。

\end{document}
